\chapter{Company Presentation and Internship Context}

\section{Company Overview}

OwnTech is an open-source power electronics initiative developed around the
idea of making power electronics more accessible, reusable, programmable, and
adaptable. Power electronics is a key technology for energy conversion: it is used
to interface photovoltaic panels, batteries, motors, electrical grids, embedded
systems, and many other energy sources or loads. However, the design of power
converters is usually complex because it combines hardware constraints,
real-time control, embedded software, measurement, safety, thermal management,
and electromagnetic compatibility.

The OwnTech logo is shown in Figure~\ref{fig:owntech_logo}.

\begin{figure}[H]
\centering

\includegraphics[width=0.38\textwidth]{OwnTech-Logo-1024x512.jpeg}
\caption{OwnTech logo.}
\label{fig:owntech_logo}
\end{figure}

The OwnTech structure combines an open-science foundation and an industrial
company. The OwnTech Foundation is placed under the aegis of the CNRS
Foundation and has existed since July 2022, while OwnTech Technologies was
created in October 2023, keeping the project open and community-oriented
while supporting industrial development. In the company presentation
material, Mr. Jean Alinei is presented as CEO of OwnTech Technologies, while
Dr. Luiz Villa is a co-founder of OwnTech. This structure illustrates the
link between engineering, research, and product development inside the
OwnTech ecosystem.

The objective of OwnTech is to reduce this complexity by providing a
complete open-source technology suite for power electronics, combining
reprogrammable power hardware, a simplified development environment,
software libraries, and monitoring tools~\cite{owntech_home,owntech_docs}.
Instead of designing a dedicated converter from zero for each application,
the user can start from a common hardware and software base and adjust it
to the required function.

This open-source approach also concerns the hardware itself. OwnTech
highlights public and editable resources such as design files, bills of
materials, assembly information, and CAD files~\cite{owntech_home}. Users
can therefore inspect the platform, reproduce experiments, and tailor the
design to specific constraints.

The project follows the principle summarized by OwnTech as ``one hardware,
many functions''. A single programmable platform can be used for several
applications, such as battery charging, photovoltaic conversion, motor
control, LED driving, or small renewable-energy
installations~\cite{owntech_home}. This versatility makes the platform
especially useful for prototyping, since the same ecosystem can be reused
to test different converter topologies and control strategies.

This accessibility objective can be understood by comparison with other
open electronics platforms. Arduino made embedded electronics accessible
to a large community of users, as illustrated in
Figure~\ref{fig:arduino_accessibility}.

\begin{figure}[H]
\centering
\includegraphics[width=0.30\textwidth]{Arduino.jpeg}
\caption{Arduino board.}
\label{fig:arduino_accessibility}
\end{figure}

In a similar way, Raspberry Pi, shown in
Figure~\ref{fig:raspberry_pi_accessibility}, made small computing platforms
accessible for education, prototyping, and embedded projects.

\begin{figure}[H]
\centering
\includegraphics[width=0.28\textwidth]{Raspberry_Pi.jpg}
\caption{Raspberry Pi board.}
\label{fig:raspberry_pi_accessibility}
\end{figure}

Micro:bit, shown in Figure~\ref{fig:microbit_accessibility}, simplified the
first contact with programming and digital systems, especially for children,
beginners, and educational activities.

\begin{figure}[H]
\centering
\includegraphics[width=0.25\textwidth]{microbit.jpeg}
\caption{Micro:bit board.}
\label{fig:microbit_accessibility}
\end{figure}

OwnTech applies the same accessibility logic to power electronics.
Figure~\ref{fig:owntech_sandbox} summarizes this broader technology sandbox
before introducing SPIN, the main hardware platform used during this internship.

\begin{figure}[H]
\centering
\includegraphics[width=1.10\textwidth]{owntech_sandbox.jpg}
\caption{OwnTech technology sandbox.}
\label{fig:owntech_sandbox}
\end{figure}

The OwnTech approach can also be described as a user-centric technology
sandbox, connecting academia, industry, makers, and fablabs to a common set
of tools, including standard hardware, an intuitive IDE, unified
documentation, and open data monitoring, so as to provide all the elements
required to create, test, share, and improve power electronics applications.

This organization creates impact by helping new ideas become real
applications in areas such as smart grids, electrical mobility, and energy
storage, and by supporting learning by doing, as students, engineers, and
makers experiment on an open platform and exchange with the community,
consistent with the OwnTech objective of giving users the tools they need to
act for the energy transition.

\section{OwnTech Boards and Tools}

The OwnTech ecosystem is a modular set of hardware boards and software tools
for power electronics, covering control, power conversion, measurement,
debugging, and monitoring~\cite{owntech_home,owntech_spin,owntech_twist}. Its
main elements are the SPIN control board, the TWIST power module, the
OwnVerter three-phase converter, and the OwnPlot monitoring interface, which
together let the user control, power, measure, and monitor converter
applications within a coherent open-source environment. This reflects a
software-defined power conversion approach, where the hardware provides the
electrical interface while the software defines the converter behaviour
through PWM generation, measurements, protections, and control algorithms.

\subsection{Main OwnTech Boards}

Table~\ref{tab:owntech_boards} summarizes the main OwnTech boards and tools
used in this project: SPIN, TWIST, OwnVerter, and
OwnPlot~\cite{owntech_twist,owntech_docs}.

\begin{table}[H]
\centering
\caption{Main OwnTech boards and tools}
\label{tab:owntech_boards}
\begin{tabularx}{\textwidth}{|p{3.0cm}|X|}
\hline
\textbf{Board or Tool} & \textbf{Main role} \\ \hline
SPIN &
OwnTech power controller board used for PWM generation, measurement acquisition,
real-time control, debugging, and communication with monitoring tools. \\ \hline
TWIST &
Programmable and stackable power module used as the power conversion stage for
DC and AC applications up to 110~V. \\ \hline
OwnVerter &
Reprogrammable 1.3~kW three-phase power converter mainly intended for
motor-control applications. \\ \hline
OwnPlot &
Monitoring interface used to interact with an application and visualize
internal variables during operation. \\ \hline
\end{tabularx}
\end{table}

With OwnTech, the user works inside a unified environment where the hardware
and software are already compatible, with a simplified graphical interface
for quick control and a development environment for adapting the application
to the required use case.

Once deployed, the platform supports data monitoring, extended documentation,
and community exchange for sharing results and solving problems. This unified
environment is illustrated in Figure~\ref{fig:owntech_user_environment}.

\begin{figure}[H]
\centering
\includegraphics[width=0.55\textwidth]{images.png}
\caption{OwnTech unified environment around the user.}
\label{fig:owntech_user_environment}
\end{figure}

For this report, the SPIN board is detailed further, as it was the central
controller used during the internship and is also a reusable control board
for several OwnTech converter platforms.

\subsection{SPIN Control Board}

OwnTech presents SPIN as a free, ergonomic, and multi-function power controller
~\cite{owntech_spin}. Its main interest is that the same control board can be
reused with different power stages, while keeping a common software environment
and development workflow. The board is shown in
Figure~\ref{fig:spin_control_board}.

Compared with a conventional microcontroller development board, SPIN provides functions that are
directly useful for switched-mode converters: PWM generation, synchronized
measurements, real-time control loops, communication with monitoring tools, and
debugging support. This makes it suitable for applications such as DC/DC
converters, inverters, motor drives, battery chargers, and laboratory
prototypes.

\begin{figure}[H]
\centering
\includegraphics[width=0.6\textwidth]{images.jpeg}
\caption{OwnTech SPIN power controller board.}
\label{fig:spin_control_board}
\end{figure}

\subsubsection{Typical SPIN Workflow}

The SPIN board is designed to reduce the number of steps required to start
controlling a power converter. The workflow remains simple: the user connects
the board through USB-C, downloads the OwnTech code and examples from GitLab,
then programs the board to drive the power stage. This lowers the entry barrier
for users who want to experiment with digital power electronics. The main
features of SPIN are summarized in Table~\ref{tab:spin_features}.


\begin{table}[H]
\centering
\caption{Main SPIN features highlighted by OwnTech}
\label{tab:spin_features}
\small
\begin{tabularx}{0.95\textwidth}{|p{2.8cm}|X|}
\hline
\textbf{Category} & \textbf{Main features} \\ \hline
Ease of use &
Native USB support, OwnPlot monitoring interface, integrated debugger, and
compatibility with MATLAB or Visual Studio environments. \\ \hline
Real-time control &
Configurable control frequency, voltage mode control, and peak current mode
control. \\ \hline
Measurement acquisition &
Six ready-to-use 12-bit ADC channels, 60~kHz acquisition rate, accurate timing
measurements, and non-blocking acquisition system. \\ \hline
PWM generation &
Eleven ready-to-use PWM signals, 55~MHz maximum frequency, 184~ps time
resolution, advanced phase-shift options, and
configurable dead-time. \\ \hline
\end{tabularx}
\end{table}

\subsubsection{Role During the Internship}

During this internship, the SPIN board was central to the first PCB project.
Its GPIO and PWM signals had to be validated through the SPOC oscilloscope
shield. The validation bench uses a SPIN board under test to generate
electrical signals, while another control stage selects which signals are
routed toward the oscilloscope. The SPOC board therefore contributes directly
to the validation of SPIN-based hardware and software behaviour.

In the second project, the same control needs appear at a higher switching
performance level: accurate PWM generation, dead-time control, isolation, and
safe switching sequences are required for the GaN half-bridge board.
\section{Internship Context}

The internship subject is also connected to a larger evolution of OwnTech
hardware. After low-voltage modular converters, the next development step is to
move toward higher voltage and higher power applications. The internship
subject targets converter modules in the range of 400--800~V and 1--10~kW,
with applications such as photovoltaic inverters, motor drives, and other
high-power energy conversion systems.

This evolution introduces new technical constraints. At higher voltage and
higher power, the design must consider insulation, creepage and clearance,
thermal limits, electromagnetic compatibility, safety, and pre-certification
requirements. The use of wide-bandgap devices such as GaN and SiC can improve
switching performance and efficiency, but it also makes PCB layout, gate
driving, parasitic effects, and measurement quality much more critical.

In this context, the internship work was organized around two different PCB
projects. Although these two boards are different, they both support the same
global objective: improving the development and validation workflow of
open-source power electronics hardware. The first project focused on automated
validation, because reliable and repeatable tests are necessary before scaling
hardware platforms. The second one concerned a GaN half-bridge board intended
for high-speed power switching experiments.

The first board is the SPIN Oscilloscope Shield, called SPOC. It is an
oscilloscope interface board used to route selected signals from a SPIN board
under test to oscilloscope inputs. It allows up to four selected GPIO or PWM
signals to be observed at the same time on a standard four-channel
oscilloscope. This board contributes to the automated test bench by making
physical signal validation more systematic and less dependent on manual
probing. The oscilloscope measurements provide feedback to the test software,
which can then decide whether the tested signal passes or fails according to
the expected electrical behaviour.

The second board is a GaN-based half-bridge board. It is intended as an
experimental platform for high-speed switching tests and future high-voltage
converter development. The board integrates digital isolation, gate-driver
stages, an isolated floating supply for the high-side driver, two 650~V GaN
transistors, and a compact DC-link path. Its design is linked to the future
high-power direction of OwnTech, where switching performance, EMI, safety, and
pre-certification become major design concerns.







\section{State of the Art: Sustainable and Modular Power Electronics}
\label{sec:sustainable_modular_power_electronics}

Power electronics is one of the enabling technologies of the energy transition.
It allows electrical energy to be converted, controlled, stored, and exchanged
between photovoltaic panels, batteries, motors, electrical grids, and embedded
systems. However, a power converter is often designed as a dedicated product
for one voltage range, one power level, one topology, and one control strategy.
When the application evolves, this type of closed design can become difficult to
inspect, repair, modify, or reuse. Sustainability in power electronics is
therefore not only a question of efficiency. It also depends on documentation,
maintenance, modularity, upgradeability, and the ability to keep hardware useful
during a longer life cycle.

The state-of-the-art discussion focuses on three
complementary concepts: open-source hardware, obsolescence management, and
software-defined modular converters. These concepts provide a useful background
for discussing how modularity, documentation, and validation can support more
sustainable power electronics hardware.

\subsection{Open-Source Hardware as a Sustainability Tool}
\label{subsec:osh_sustainability_tool}

Open-source hardware gives users access to the design information of a physical
system. In practice, this means that the hardware can be operated, studied,
modified, reproduced, manufactured, and sold under the conditions of the
associated licence. Bonvoisin et al. explain that open-source hardware requires
standardised documentation practices in order to make these rights usable in
practice~\cite{bonvoisin2020standardisation}. A schematic or a PCB file alone is
not enough: users also need interface descriptions, assembly information,
manufacturing data, software, tests, safety limits, and maintenance-oriented
information. Figure~\ref{fig:open_source_hardware_rights} illustrates the
main rights associated with open-source hardware.

\begin{figure}[htbp]
    \centering
    \IfFileExists{images/chapitre 1/slide_10.png}{%
        \includegraphics[width=0.80\textwidth]{images/chapitre 1/slide_10.png}%
    }{%
        \fbox{\parbox[c][4.2cm][c]{0.85\textwidth}{\centering Missing figure file}}%
    }
    \caption{Open-source hardware.}
    \label{fig:open_source_hardware_rights}
\end{figure}

For power electronics, this openness has a direct technical value. A converter
combines power semiconductors, gate drivers, sensing circuits, protection
functions, embedded control, communication, thermal design, and electromagnetic
compatibility constraints. If these elements are documented, the converter is
not only a product to be used; it also becomes a technical object that can be
understood, validated, repaired, and adapted by engineers, researchers,
students, and makers.

\subsection{Obsolescence in Power Electronics}
\label{subsec:obsolescence_power_electronics}

Obsolescence describes the situation in which a product, component, or system
becomes outdated, unavailable, unsuitable, or less useful with respect to current
needs. In industrial systems, obsolescence management is also treated as a
life-cycle activity intended to anticipate and reduce the impact of unavailable
or outdated items. Mellal identifies several forms of
obsolescence, including technical, functional, planned, psychological, and
optional obsolescence~\cite{mellal2020obsolescence}. In power electronics,
technical and functional obsolescence are the most critical for industrial and
research platforms.

\begin{figure}[htbp]
    \centering
    \IfFileExists{images/chapitre 1/slide_12.png}{%
        \includegraphics[width=0.80\textwidth]{images/chapitre 1/slide_12.png}%
    }{%
        \fbox{\parbox[c][4.2cm][c]{0.80\textwidth}{\centering Missing figure file}}%
    }
    \caption{Main obsolescence mechanisms in power electronics.}
    \label{fig:power_electronics_obsolescence}
\end{figure}

Figure~\ref{fig:power_electronics_obsolescence} summarizes the main forms of
obsolescence that can affect a power electronics system. Technical obsolescence
comes from component unavailability, reference discontinuation, or supplier
disappearance. Functional obsolescence appears when the converter still works,
but no longer matches the required power level, energy flow, or conversion type.
Planned obsolescence is mainly related to sizing choices made during design, for
example with electrolytic capacitors whose lifetime depends strongly on
temperature and electrical stress.

The figure also includes psychological and optional obsolescence. Psychological
obsolescence is linked to footprint, weight, or style, and is mostly a B2C
criterion, although it can also influence the perceived usability of a platform.
Optional obsolescence occurs when a converter cannot adapt to changes in its
environment, for example because its interfaces, sensing possibilities, control
software, or hardware architecture are too fixed. This is important for OwnTech,
whose objective is to build reusable converter platforms that can be maintained
and adapted instead of replaced as black boxes.

\subsection{Modularity and Software-Defined Power Conversion}
\label{subsec:modularity_sdpc}

Modularity is one way to reduce the impact of obsolescence: the converter can
be separated into functional blocks with clear interfaces, such as control,
sensing, communication, isolation, gate driving, power switching, and load
interface, so that an obsolete or defective block can be redesigned or
replaced without changing the whole platform.

This idea is related to the PEBB concept, where converters are organised as
reusable bricks combining switching devices, drivers, sensing, and local
control, which can then be assembled into a PCA to scale voltage, current, or
power. OwnTech connects this logic to SDPC, where the hardware stays generic
while the software defines the control loops, protections, and effective
conversion function. Figure~\ref{fig:pebb_pca_sdpc_concepts} summarizes these
related concepts.

\begin{figure}[htbp]
    \centering
    \IfFileExists{images/chapitre 1/slide_16.png}{%
        \includegraphics[width=0.70\textwidth]{images/chapitre 1/slide_16.png}%
    }{%
        \fbox{\parbox[c][4.2cm][c]{0.8\textwidth}{\centering Missing figure file}}%
    }
    \caption{PEBB, PCA, and SDPC concepts.}
    \label{fig:pebb_pca_sdpc_concepts}
\end{figure}

Rahmani et al. discuss the opportunities of modularity for the eco-design of
power converters, especially through repair, reuse, and replacement of
sub-parts~\cite{rahmani2021modularity}. This is particularly relevant in power
electronics because the stressed elements of a converter do not age at the same
speed. Power semiconductors, gate drivers, DC-link capacitors, filters, thermal
interfaces, control electronics, and communication circuits may all have
different life-cycle constraints.

In the context of this internship, this modularity principle is directly linked
to the separation between the control and communication functions, the
instrumentation functions, the isolation stage, the gate-driver stage, and the
GaN half-bridge power stage. This separation makes the system easier to analyse
and validate, and it also makes future upgrades more practical. For example, a
new power stage could be developed to test another GaN or SiC device while
keeping the same control and instrumentation environment. Conversely, the
control or communication board could evolve without requiring a complete redesign
of the power stage.

\subsection{Documentation and Maintenance}
\label{subsec:documentation_maintenance}

A major maintenance issue in power electronics is that converters are often treated as black boxes. When a fault appears, the complete
unit is frequently replaced because the internal design, diagnostic information,
and repair procedure are not available. Open-source hardware can help reduce
this problem only if the documentation covers the full life cycle of the
product, not only its manufacturing phase. Figure~\ref{fig:repairable_converter_architecture} illustrates Modular architecture for a repairable power converter.

\begin{figure}[htbp]
    \centering
    \IfFileExists{images/chapitre 1/slide_19.png}{%
        \includegraphics[width=0.66\textwidth]{images/chapitre 1/slide_19.png}%
    }{%
        \fbox{\parbox[c][4.2cm][c]{0.9\textwidth}{\centering Missing figure file}}%
    }
    \caption{Modular architecture for a repairable power converter.}
    \label{fig:repairable_converter_architecture}
\end{figure}

This representation is directly related to the modular approach used in the
internship work. A converter can be viewed as an assembly of functional blocks:
control, instrumentation, gate driving, power semiconductors, filtering, and
load interface. If one block becomes defective or obsolete, the objective is to
replace or redesign only this block instead of replacing the complete converter.
This approach supports repairability, upgradeability, and reduction of material
waste, provided that the interfaces, diagnostic points, replacement options, and
maintenance information of each block are documented clearly.

\subsection{Link with the Internship Work}
\label{subsec:internship_positioning_sustainability}

This link with the state of the art is mainly visible through the GaN
half-bridge board. The board is intended as a reusable experimental power stage
for high-speed and high-voltage switching tests. It follows the modular logic
described above by separating the control board, measurement interface,
isolation, gate-driver stage, and power stage.

This separation makes the system easier to understand, validate, modify, and
upgrade. Instead of redesigning one large board, several parameters can be
adjusted by changing only the block or small board concerned. For example, the
power stage can evolve toward another GaN or SiC device while keeping the same
control and measurement environment. Its power capability can also be adapted by
replacing the selected MOSFET with a more powerful device, or by studying
series-connected GaN MOSFETs when higher voltage or power levels are required.
In this sense, the internship contributes to the OwnTech objective of building
power electronics hardware that is more open, modular, maintainable, and
adaptable.
